\customchapter{RESUMEN}
% Correccion conjunta resumen/abstract r02. Trazabilidad: redaccion/rondas/resumen-r02-respuesta.md.

Los implantes metálicos producen rayas y zonas oscuras en la tomografía computarizada pélvica que pueden ocultar los límites del hueso. Como hay pocas tomografías con implantes y anotaciones óseas, se plantea generar imágenes con implantes a partir de tomografías sin metal. Este trabajo propuso una cadena que coloca un tornillo transilíaco-transsacro y sintetiza su artefacto metálico local en tomografías sin implante del conjunto CTPelvic1K; sus dos componentes se evaluaron por separado.

Primero se comprobó si una red que comprime y reconstruye imágenes conservaba los valores del hueso necesarios para sintetizar en un espacio comprimido. Luego, un muestreador propuso posiciones tridimensionales del tornillo y se definió una medida de admisibilidad quirúrgica basada en la viabilidad del corredor óseo y el grado de brecha cortical, que registra cuánto sobresale el tornillo del hueso. Para generar la imagen se empleó difusión directamente sobre sus valores, con codificación multiventana; solo se regeneraron el implante y una banda de unos 12~mm a su alrededor. La intensidad de las rayas generadas, medida mediante la \emph{streak amplitude}, se comparó con la de un protocolo de simulación física y con la inserción del tornillo sin artefacto como control.

La prueba de compresión falló: ninguna de las seis configuraciones evaluadas alcanzó el límite fijado de 25 unidades Hounsfield de error óseo; el menor error medio fue de 61.72 unidades Hounsfield en 34 pacientes. Por ese resultado, la síntesis se implementó directamente sobre los valores de la imagen. En 72 pacientes se generaron 3\,600 posiciones; la distribución de sus grados de brecha quedó a distancias de 0.206 y 0.230 grados de las dos series clínicas de referencia, separadas por 0.395 grados. Esa cercanía es descriptiva, pues no se había fijado un criterio para declarar satisfactoria la colocación. En 13 pacientes, la intensidad media de las rayas sintetizadas fue 301.0 unidades Hounsfield menor que la del protocolo físico con titanio. El intervalo de confianza del 90\,\% para esa diferencia cruzó el límite de equivalencia: no se pudo establecer que ambas intensidades fueran equivalentes.

La cadena quedó implementada, pero no se estableció la equivalencia de la intensidad de sus rayas con la del protocolo físico. La autora juzgó insuficiente su realismo visual; además, el artefacto sintetizado no reprodujo la caída de intensidad de las rayas al alejarse del metal. No se evaluó si estas imágenes mejoran el entrenamiento de redes de segmentación ósea.

\noindent \textbf{Palabras clave:}\\
\noindent Síntesis de imágenes médicas; Artefacto metálico; Tomografía computarizada pélvica; Modelos de difusión; Fijación iliosacra
